\newcommand{\hupar}[1]{\vspace{12pt} \noindent \textbf{#1}\quad}
\begin{abstract}
We proposed \textbf{P}recomputed \textbf{R}adiance \textbf{T}ransfer of \textbf{G}aussian \textbf{S}plats (PRTGS), a real-time high-quality relighting method for Gaussian splats in low-frequency lighting environments that captures soft shadows and interreflections by precomputing 3D Gaussian splats' radiance transfer. Existing studies have demonstrated that 3D Gaussian splatting (3DGS) outperforms neural fields' 
 efficiency for dynamic lighting scenarios. However, the current relighting method based on 3DGS is still struggling to compute high-quality shadow and indirect illumination in real time for dynamic light, leading to unrealistic rendering results. We solve this problem by precomputing the expensive transport simulations required for complex transfer functions like shadowing, the resulting transfer functions are represented as dense sets of vectors or matrices for every Gaussian splat. We introduce distinct precomputing methods tailored for training and rendering stages, along with unique ray tracing and indirect lighting precomputation techniques for 3D Gaussian splats to accelerate training speed and compute accurate indirect lighting related to environment light. Experimental analyses demonstrate that our approach achieves state-of-the-art visual quality while maintaining competitive training times and allows high-quality real-time (30+ fps) relighting for dynamic light and relatively complex scenes at 1080p resolution.
%As a neuromorphic sensor with high temporal resolution, the spike camera has great potential in high-speed vision applications. However, in low-light conditions, the scarcity of information severely limits the quality of collected spike streams. To address it, we first propose a deep learning framework that transforms low-light spike streams into enhanced spike streams, with the spike firing frequency serving as a bridge. LS2ES first uses a network to learn to enhance the spike firing frequency of the low-light spike stream. Then, a decoder is introduced to generate the enhanced spike stream based on the enhanced spike firing frequency sequences. LS2ES effectively improves the spatial-temporal distribution of spikes in low-light spike streams, leading to more concentrated information. 
%In order to evaluate its performance on low-light spike streams, we construct a low-light spike streams dataset for driving scenes using the latest generation of spike cameras. By reconstructing both the low-light spike stream and the enhanced spike stream separately, we find that LS2ES can improve the brightness and details of scenes. Furthermore, the enhanced spike stream potentially benefits other spike-based tasks, e.g., super-resolution and optical flow estimation.
\end{abstract}

%作为一个有着高时间分辨率的神经形态传感器，脉冲相机在高速视觉应用上有着巨大的潜力。然而, 在低光场景中，匮乏的信息将严重地限制被采集的脉冲流质量。为此, 我们首次提出了一个从低光脉冲流到增强脉冲流的深度学习框架where脉冲发放频率作为一个桥梁。LS2ES首先使用一个网络去学习增强低光脉冲流的脉冲发放频率。然后一个解码器被引入去生成增强脉冲流根据增强脉冲发放频率序列。LS2ES能有效地改善低光脉冲流中脉冲的时空分布使信息更加聚集
%
%。为了评估它在低光脉冲流上的表现，我们首次使用最新一代的相机构建了a low-light spike streams dataset for driving scenarios. 通过分别重构低光脉冲流与增强脉冲流，我们发现LS2ES有效地改善了场景的亮度与细节。此外，增强脉冲流还潜在地对其他脉冲域任务有着受益, e.g.,。
%
%它能有效地改善低光脉冲流中脉冲的时空分布使信息更加聚集。
